% ID: SEL-D0014-BB-E04
% ESTADO: validada. S015; tres pasadas del mismo autor, sin segundo revisor.
\textbf{a) Continuidad y derivabilidad en cero.}\par\nopagebreak[4]
El dominio es $\mathbb R\setminus\{1\}$. En $0$, ambas expresiones tienen límite $1$ y $f(0)=1$; es continua. Las derivadas de las ramas son $2x+1$ y $(1-x)^{-2}$. Sus valores laterales en $0$ son ambos $1$:
\[
\boxed{f\text{ es continua y derivable en }0,\qquad f'(0)=1.}
\]
\textbf{b) Monotonía y curvatura.}\par\nopagebreak[4]
En la rama izquierda $f'=2x+1$ cambia de negativa a positiva en $x=-1/2$. En la derecha, $f'=1/(1-x)^2>0$ en cada intervalo de su dominio. Así, decrece en $(-\infty,-1/2]$ y crece en $[-1/2,1)$ y en $(1,+\infty)$, por separado. No se unen intervalos a través de la asíntota. Hay mínimo relativo en $x=-1/2$, de valor $3/4$.
Las segundas derivadas son
\[
f''(x)=\begin{cases}2&x<0,\\2/(1-x)^3&x>0,\ x\ne1.\end{cases}
\]
Además $f''(0)=2$, pues las derivadas laterales de $f'$ coinciden. Por tanto, $f''>0$ en $(-\infty,1)$: \textbf{curvatura hacia arriba}; y $f''<0$ en $(1,+\infty)$: \textbf{curvatura hacia abajo}. No hay puntos de inflexión: en $0$ no cambia el signo y $1$ no pertenece al dominio.

\textbf{c) Asíntotas.}\par\nopagebreak[4]
\[
\lim_{x\to1^-}f(x)=+\infty,\qquad\lim_{x\to1^+}f(x)=-\infty,
\]
luego \boxed{x=1} es vertical. Como $\lim_{x\to+\infty}1/(1-x)=0$, \boxed{y=0} es horizontal hacia $+\infty$. A la izquierda la expresión cuadrática no admite asíntota horizontal ni oblicua: $f(x)/x=x+1+1/x$ no tiene límite finito. No hay otras asíntotas rectilíneas.
